\documentclass[11pt,letterpaper]{article}

\usepackage[utf8]{inputenc}
\usepackage[T1]{fontenc}
\usepackage[spanish,es-noshorthands]{babel}
\usepackage{spectral}
\usepackage[margin=2.5cm,headheight=48pt,headsep=18pt]{geometry}
\usepackage{fancyhdr}
\usepackage{graphicx}
\usepackage{xcolor}
\usepackage{listings}
\usepackage{tikz}
\usepackage{booktabs}
\usepackage{enumitem}
\usepackage{float}
\usepackage[hidelinks]{hyperref}
\usetikzlibrary{arrows.meta,positioning}

\newcommand{\curso}{CC4303-1 Redes}
\newcommand{\tituloinforme}{Actividad: Sockets orientados a conexión con Stop \& Wait}

\renewcommand{\sectionmark}[1]{\markboth{#1}{}}
\setlength{\parskip}{0.6em}
\setlength{\parindent}{0pt}

% Cabecera y pie de página con el mismo estilo del informe de práctica.
\pagestyle{fancy}
\fancyhf{}
\fancyhead[L]{\small\leftmark}
\fancyfoot[L]{\small Informe \tituloinforme}
\fancyfoot[R]{\small CC4303-1}
\renewcommand{\headrulewidth}{0.4pt}
\renewcommand{\footrulewidth}{0.4pt}

\fancypagestyle{portada}{%
  \fancyhf{}
  \fancyhead[L]{\small\begin{tabular}[b]{@{}l@{}}
    Universidad de Chile\\
    Facultad de Ciencias Físicas y Matemáticas\\
    Departamento de Ciencias de la Computación\\
    \curso
  \end{tabular}}
  \fancyhead[R]{\includegraphics[height=1.3cm]{logo_fcfm.png}}
  \renewcommand{\headrulewidth}{0.4pt}
  \renewcommand{\footrulewidth}{0pt}
}

\definecolor{fondocodigo}{gray}{0.96}
\lstset{
  basicstyle=\ttfamily\footnotesize,
  backgroundcolor=\color{fondocodigo},
  frame=single,
  rulecolor=\color{gray!50},
  breaklines=true,
  columns=fullflexible,
  keepspaces=true,
  showstringspaces=false,
  literate={á}{{\'a}}1 {é}{{\'e}}1 {í}{{\'i}}1 {ó}{{\'o}}1 {ú}{{\'u}}1
           {Á}{{\'A}}1 {ñ}{{\~n}}1
}

% Diagramas de secuencia: dos líneas de tiempo verticales.
\tikzset{
  msg/.style={-{Stealth[length=2.5mm]},thick},
  lost/.style={-{Stealth[length=2.5mm]},thick,dashed,red!70!black},
  note/.style={font=\footnotesize,align=left},
  proc/.style={font=\footnotesize\itshape,text=blue!50!black},
}

\begin{document}

\begin{titlepage}
  \thispagestyle{portada}
  \centering
  \vspace*{5cm}
  {\LARGE Informe\par}
  \vspace{0.3cm}
  {\Large \tituloinforme\par}
  \vspace{0.4cm}
  {\large Juan Pablo Bustamante Flaño\par}
  \vfill
  Correo electrónico: juan.bustamante.f@ug.uchile.cl\\
  Curso: \curso{} --- Primavera 2026\\
  Profesora: Ivana Bachmann\\
  Módulo 3: Protocolos de transporte, semanas 6-8\\
  Fecha: 19/10/2026
  \vspace*{3cm}
\end{titlepage}

\section{Introducción}

En esta actividad se implementó una versión simplificada de TCP sobre sockets UDP,
usando Stop \& Wait como mecanismo de control de errores. El resultado es la clase
\texttt{SocketTCP}, que ofrece una interfaz parecida a la de los sockets orientados a
conexión de Python (\texttt{bind}, \texttt{connect}, \texttt{accept}, \texttt{send},
\texttt{recv}, \texttt{close} y \texttt{recv\_close}), junto a un cliente y un servidor
que la utilizan para transmitir un archivo de forma íntegra, incluso cuando se pierden
segmentos en ambas direcciones.

\begin{center}
\begin{tabular}{@{}ll@{}}
  \toprule
  Archivo & Contenido \\
  \midrule
  \texttt{SocketTCP.py} & Clase \texttt{SocketTCP} y constantes de configuración. \\
  \texttt{cliente.py}, \texttt{servidor.py} & Cliente y servidor de entrega. \\
  \texttt{cliente\_udp.py}, \texttt{servidor\_udp.py} & Programas provisorios de los pasos 1 y 3. \\
  \texttt{test\_cliente.py}, \texttt{test\_servidor.py} & Pruebas de \texttt{send}, \texttt{recv} y cierre. \\
  \texttt{archivo.txt} & Archivo de prueba. \\
  \bottomrule
\end{tabular}
\end{center}

Uso de los programas de entrega:

\begin{lstlisting}
% python3 servidor.py localhost 8000 > recibido.txt
% python3 cliente.py localhost 8000 < archivo.txt
\end{lstlisting}

El servidor escribe el archivo recibido en salida estándar y los mensajes de debug en
salida de error, de modo que ambos flujos no se mezclan y es posible comparar el archivo
recibido con el original usando \texttt{cmp}.

\section{Desarrollo}

\subsection{Paso 1: cliente y servidor provisorios}

\texttt{cliente\_udp.py} pide la ruta del archivo con \texttt{input()}, lo carga completo
en memoria y lo envía a \texttt{(localhost, 8000)} en trozos de a lo más 16 bytes.
\texttt{servidor\_udp.py} recibe los trozos e imprime su contenido. Como todavía no se
conoce el largo del mensaje, el servidor se queda esperando más mensajes una vez que
recibió todo.

\textbf{Test 1.} Al enviar \texttt{archivo.txt} (304 bytes, 19 trozos) sin pérdidas, el
servidor imprimió el archivo completo y sin errores.

\textbf{Test 2.} Se indujo pérdida con \texttt{netem}, con la configuración del
enunciado:

\begin{lstlisting}
% sudo tc qdisc add dev lo root netem loss 20.0% delay 0.5s
% echo archivo.txt | python3 cliente_udp.py
\end{lstlisting}

El servidor imprimió un texto con fragmentos faltantes (por ejemplo, desde
\texttt{incl} salta a \texttt{erden segmentos}):

\begin{lstlisting}
Wait.
Este archivo tiene bastante más de 16 bytes, por lo que el cliente
debe dividirlo en varios segmentos y el servidor debe rearmarlo
de forma íntegra inclerden segmentos o ACKs.
Áéíóú ñ: también se prueban caracteres UTF-8 de más de un byte.
\end{lstlisting}

Esto muestra que UDP por sí solo no garantiza la entrega y motiva el resto de la
actividad.

\subsection{Paso 2: la clase \texttt{SocketTCP}}

El constructor no recibe parámetros y crea un socket UDP (\texttt{AF\_INET},
\texttt{SOCK\_DGRAM}). Además inicializa el estado que se usa durante la conexión: las
direcciones local y de destino, los contadores \texttt{seq} y \texttt{expected\_seq}, el
\texttt{seq} del último \texttt{ACK} del handshake, la cantidad de bytes del mensaje que
quedan por entregar (\texttt{message\_pending}), un buffer con los bytes recibidos que no
cupieron en el \texttt{buff\_size} pedido (\texttt{recv\_buffer}), una lista de segmentos
recibidos antes de tiempo (\texttt{pending\_segments}) y la lista de conexiones ya
aceptadas.

\subsection{Paso 3: headers TCP}

Se implementaron \texttt{parse\_segment} y \texttt{create\_segment}, que convierten entre
bytes y un diccionario con las llaves \texttt{syn}, \texttt{ack}, \texttt{fin},
\texttt{seq} y \texttt{data}. \textbf{Decisión de diseño:} el header es texto con el
formato \texttt{SYN|ACK|FIN|SEQ|}, donde cada flag es \texttt{0} o \texttt{1} y
\texttt{SEQ} usa siempre 10 dígitos; por ejemplo, \texttt{b'1|0|0|0000000042|hola'}.
Así todos los headers miden 17 bytes, los datos empiezan siempre en la posición 17
(pueden contener \texttt{|} sin confundir al parser) y el socket UDP interno usa un buffer
fijo de $17 + 16 = 33$ bytes. El número de secuencia avanza en los bytes de datos
enviados y cada \texttt{ACK} lleva el siguiente byte esperado, como en el video.

Al comienzo de \texttt{SocketTCP.py} se definen como variables globales el timeout
(\texttt{TIMEOUT = 1.5} s, mayor que el RTT de 1 s que produce \texttt{netem delay
0.5s}) y el modo debug (\texttt{DEBUG = True}, que imprime en \texttt{stderr} cada
segmento, timeout y retransmisión).

Para el test, el cliente y el servidor provisorios aceptan la opción
\texttt{-{}-headers}: el cliente encapsula cada trozo con \texttt{create\_segment} y el
servidor extrae los datos con \texttt{parse\_segment}. El servidor rearmó el archivo
completo, igual que en el test 1.

\subsection{Paso 4: 3-way handshake}

\begin{itemize}[nosep]
  \item \texttt{bind(address)} asocia el socket UDP a la dirección y la guarda.
  \item \texttt{connect(address)} elige un número de secuencia inicial $x$ aleatorio
        entre 0 y 100, envía \texttt{SYN seq=x} y espera \texttt{SYN+ACK seq=x+1}.
        Luego envía \texttt{ACK seq=x+2} a la dirección desde donde llegó el
        \texttt{SYN+ACK}, que es la del nuevo socket del servidor.
  \item \texttt{accept()} espera un \texttt{SYN} y crea un nuevo \texttt{SocketTCP}
        asociado a la misma IP, en el primer puerto libre después del que escucha. Prueba
        con \texttt{bind} el puerto siguiente y, si está ocupado (\texttt{OSError}), sigue
        con el próximo. Desde ese socket envía el \texttt{SYN+ACK} y espera el
        \texttt{ACK}. Retorna el nuevo socket y su dirección, que es distinta a la del
        socket que escucha.
\end{itemize}

Al terminar el handshake ambos extremos tienen \texttt{seq = expected\_seq = x+2}. La
figura~\ref{fig:handshake} muestra el diagrama y qué función ejecuta cada paso.

\begin{figure}[H]
\centering
\begin{tikzpicture}[yscale=0.9]
  \node[font=\bfseries] at (0,0.6) {Cliente};
  \node[font=\bfseries] at (8,0.6) {Servidor};
  \draw[thick] (0,0) -- (0,-6.2);
  \draw[thick] (8,0) -- (8,-6.2);
  \draw[msg] (0,-0.5) -- node[above,sloped] {\texttt{SYN seq=x}} (8,-1.7);
  \draw[msg] (8,-2.4) -- node[above,sloped] {\texttt{SYN+ACK seq=x+1}} (0,-3.6);
  \draw[msg] (0,-4.3) -- node[above,sloped] {\texttt{ACK seq=x+2}} (8,-5.5);
  \node[proc,anchor=east,align=right] at (-0.4,-3.1) {\texttt{connect(address)}\\$x$ aleatorio en $[0,100]$};
  \node[proc,anchor=west,align=left] at (8.3,-1.7) {\texttt{accept()}: socket que\\escucha recibe el \texttt{SYN}};
  \node[proc,anchor=west,align=left] at (8.3,-3.4) {\texttt{accept()}: el nuevo socket\\envía el \texttt{SYN+ACK}};
  \node[proc,anchor=west,align=left] at (8.3,-5.5) {\texttt{accept()} retorna\\(nuevo socket, dirección)};
  \node[proc,anchor=east,align=right] at (-0.4,-5.6) {\texttt{connect} retorna};
\end{tikzpicture}
\caption{3-way handshake y su relación con \texttt{connect} y \texttt{accept}.}
\label{fig:handshake}
\end{figure}

\textbf{Test.} Con el código del enunciado para el cliente y el servidor, el handshake
terminó con éxito:

\begin{lstlisting}
(debug) Enviando [SYN seq=20] a ('localhost', 8000)
(debug) Recibido [SYN+ACK seq=21] desde ('127.0.0.1', 8001), enviando ACK seq=22
\end{lstlisting}

\subsection{Paso 5: Stop \& Wait con \texttt{send} y \texttt{recv}}

\textbf{\texttt{send(message)}.} Primero envía un segmento con el largo del mensaje,
codificado como texto decimal (por ejemplo, \texttt{b'304'}), y luego el mensaje en trozos
de 16 bytes. Cada segmento se envía con Stop \& Wait: se envía, se espera el \texttt{ACK}
con \texttt{seq} igual a \texttt{seq + len(datos)} y, si se cumple el timeout, se
retransmite. El timeout se implementa como en el material del curso: se configura con
\texttt{settimeout(TIMEOUT)} y el error \texttt{socket.timeout} que lanza
\texttt{recvfrom} se captura con \texttt{try}/\texttt{except} para retransmitir. Los
segmentos que llegan y no son el \texttt{ACK} esperado se ignoran y se sigue esperando. Al recibir el \texttt{ACK}, \texttt{seq} avanza y se envía el siguiente
segmento. El primer segmento usa el último \texttt{seq} guardado, así que varios
\texttt{send} seguidos continúan la numeración.

\textbf{\texttt{recv(buff\_size)}.} Usa dos atributos: \texttt{message\_pending}, los
bytes del mensaje actual que aún no se entregan, y \texttt{recv\_buffer}, los bytes ya
recibidos y confirmados que todavía no se entregan. El algoritmo es:

\begin{enumerate}[nosep]
  \item Si \texttt{message\_pending == 0}, comienza un mensaje nuevo: el primer segmento
        trae el largo, que se guarda en \texttt{message\_pending}.
  \item Se calcula \texttt{objetivo = min(message\_pending, buff\_size)}.
  \item Se reciben segmentos en orden y se agregan a \texttt{recv\_buffer} hasta que este
        tenga al menos \texttt{objetivo} bytes.
  \item Se retornan los primeros \texttt{objetivo} bytes. El resto queda en
        \texttt{recv\_buffer} y \texttt{message\_pending} disminuye en \texttt{objetivo}.
\end{enumerate}

\textbf{Decisión de diseño: bytes sobrantes.} Como los segmentos traen hasta 16 bytes,
puede ocurrir que \texttt{len(recibido) > buff\_size}, por ejemplo con
\texttt{buff\_size = 14}. En ese caso los bytes que sobran no se descartan ni se dejan
sin confirmar: ya se confirmaron con su \texttt{ACK} y quedan en \texttt{recv\_buffer}
para la siguiente llamada a \texttt{recv}. Así \texttt{recv} nunca retorna más de
\texttt{buff\_size} bytes y no se pierden datos. Con esto, \texttt{recv} retorna apenas se
cumple \texttt{len(recibido) >= min(message\_length, buff\_size)}, como pide el enunciado.

\textbf{Recepción de cada segmento.} El receptor solo acepta un segmento de datos si su
\texttt{seq} es igual a \texttt{expected\_seq}. En ese caso avanza
\texttt{expected\_seq} y envía \texttt{ACK seq=expected\_seq}. Si llega un segmento con un
\texttt{seq} anterior, es una retransmisión porque se perdió el \texttt{ACK}: se descartan
sus datos y se reenvía el \texttt{ACK} con el \texttt{expected\_seq} actual.

\textbf{Test.} Sin pérdidas, los tests del enunciado pasaron:

\begin{lstlisting}
Test 1 received: b'Mensje de len=16'
Test 1: Passed
Test 2 received: b'Mensaje de largo 19'
Test 2: Passed
Test 3 received: b'Mensaje de largo 19'
Test 3: Passed
\end{lstlisting}

\subsection{Paso 6: cierre de conexión}

\texttt{close()} implementa el lado del Host A: envía \texttt{FIN seq=s}, espera
\texttt{FIN+ACK seq=s+1} y responde \texttt{ACK seq=s+2}. \texttt{recv\_close()}
implementa el lado del Host B: espera el \texttt{FIN}, responde \texttt{FIN+ACK} y espera
el último \texttt{ACK}. En ambos casos se cierra el socket UDP al terminar. Se decidió
combinar el \texttt{ACK} del \texttt{FIN} y el \texttt{FIN} de B en un solo segmento
\texttt{FIN+ACK}, que es lo que espera \texttt{close()} en el paso 8.

\textbf{Test.} Al agregar \texttt{close()} al final del cliente de prueba y
\texttt{recv\_close()} al final del servidor, la conexión se cerró con éxito:

\begin{lstlisting}
# cliente                                   # servidor
(debug) Enviando [FIN seq=126]               (debug) Recibido [FIN seq=126]
(debug) Recibido [FIN+ACK seq=127],          (debug) Enviando [FIN+ACK seq=127]
        enviando [ACK seq=128] 3 veces       (debug) Recibido [ACK seq=128],
(debug) Conexión cerrada                             conexión cerrada
\end{lstlisting}

\subsection{Paso 7: handshake con pérdidas}

\texttt{connect} y \texttt{accept} se modificaron para usar Stop \& Wait:
\texttt{connect} retransmite el \texttt{SYN} si no llega el \texttt{SYN+ACK} antes del
timeout, y \texttt{accept} retransmite el \texttt{SYN+ACK} si no llega el \texttt{ACK}.
Como el cliente retransmite el \texttt{SYN} al socket que escucha, el servidor guarda las
conexiones ya aceptadas (dirección del cliente, \texttt{seq} inicial) e ignora los
\texttt{SYN} repetidos de esas conexiones en llamadas posteriores a \texttt{accept}.

\textbf{Caso borde: se pierde el último \texttt{ACK} del handshake.} El cliente no tiene
forma de saber que el \texttt{ACK} se perdió: \texttt{connect} ya retornó y el cliente
pasa a \texttt{send}. El servidor, en cambio, sigue en \texttt{accept} esperando el
\texttt{ACK}, retransmite el \texttt{SYN+ACK} y no confirma los datos del cliente. La
figura~\ref{fig:borde} muestra este caso y la solución.

\begin{figure}[H]
\centering
\begin{tikzpicture}[yscale=0.9]
  \node[font=\bfseries] at (0,0.6) {Cliente};
  \node[font=\bfseries] at (8,0.6) {Servidor};
  \draw[thick] (0,0) -- (0,-10.4);
  \draw[thick] (8,0) -- (8,-10.4);
  \draw[msg] (0,-0.4) -- node[above,sloped] {\texttt{SYN seq=x}} (8,-1.4);
  \draw[msg] (8,-1.8) -- node[above,sloped] {\texttt{SYN+ACK seq=x+1}} (0,-2.8);
  \draw[lost] (0,-3.2) -- node[above,sloped] {\texttt{ACK seq=x+2}} (4.5,-3.8);
  \node[red!70!black,font=\Large] at (4.7,-3.85) {$\times$};
  \node[proc,anchor=east,align=right] at (-0.3,-3.2) {\texttt{connect} retorna};
  \node[proc,anchor=east,align=right] at (-0.3,-4.4) {\texttt{send}};
  \draw[msg] (0,-4.4) -- node[above,sloped] {\texttt{DATA seq=x+2 (largo)}} (8,-5.4);
  \node[proc,anchor=west,align=left] at (8.3,-5.4) {\texttt{accept}: segmento sin\\\texttt{SYN} con \texttt{seq=x+2}\\$\Rightarrow$ ACK implícito};
  \node[proc,anchor=west,align=left] at (8.3,-7.0) {se guarda en\\\texttt{pending\_segments}};
  \node[proc,anchor=west,align=left] at (8.3,-8.2) {\texttt{recv} lo procesa};
  \draw[msg] (8,-8.2) -- node[above,sloped] {\texttt{ACK seq=x+2+L}} (0,-9.2);
  \draw[densely dotted] (-0.2,-6.2) -- (8.2,-6.2);
  \node[note,fill=white] at (4,-6.2) {Variante: el servidor alcanza a retransmitir};
  \draw[msg] (8,-6.6) -- node[above,sloped] {\texttt{SYN+ACK seq=x+1}} (0,-7.6);
  \node[proc,anchor=east,align=right] at (-0.3,-7.6) {\texttt{send}/\texttt{recv}:\\reenvía \texttt{ACK seq=x+2}};
\end{tikzpicture}
\caption{Caso borde: se pierde el último \texttt{ACK} del handshake.}
\label{fig:borde}
\end{figure}

\textbf{Solución.} Se usan dos mecanismos complementarios:

\begin{enumerate}[nosep]
  \item \textbf{ACK implícito en \texttt{accept}.} Mientras espera el \texttt{ACK}, el
        servidor también acepta cualquier segmento sin \texttt{SYN} cuyo \texttt{seq} sea
        $x+2$. Ese número solo lo puede usar el cliente después de recibir el
        \texttt{SYN+ACK}, así que confirma que el handshake terminó. Si el segmento trae
        datos (o es un \texttt{FIN}), no se descarta: se guarda en
        \texttt{pending\_segments} del nuevo socket y es lo primero que procesa
        \texttt{recv} (o \texttt{recv\_close}), que lo confirma con su \texttt{ACK}.
  \item \textbf{Reenvío del \texttt{ACK} en \texttt{send} y \texttt{recv}.} Si el
        cliente recibe un \texttt{SYN+ACK} repetido mientras espera un \texttt{ACK} de
        datos (o mientras espera datos), reenvía el \texttt{ACK seq=x+2} del handshake.
        Esto cubre el caso en que el cliente llama a \texttt{recv} en vez de
        \texttt{send} y nunca enviaría un segmento que sirva como ACK implícito.
\end{enumerate}

\textbf{Test.} En las pruebas con \texttt{netem} (sección~\ref{sec:netem}) el ACK
implícito se usó en 4 de 30 ejecuciones y el reenvío del \texttt{ACK} ante un
\texttt{SYN+ACK} repetido en 3. En la primera de estas ejecuciones, el \texttt{ACK} final
del handshake no llegó y el cliente pasó directamente a \texttt{send}. El servidor, que
seguía en \texttt{accept}, completó el handshake al recibir el primer segmento de datos
(\texttt{seq=12}, que es $x+2$):

\begin{lstlisting}
# cliente
(debug) Enviando [SYN seq=10] a ('localhost', 9265)
(debug) Recibido [SYN+ACK seq=11] desde ('127.0.0.1', 9266), enviando ACK seq=12
(debug) Enviando [DATA seq=12 data=b'304']
# servidor
(debug) Recibido [SYN seq=10] desde ('127.0.0.1', 58672)
(debug) Enviando [SYN+ACK seq=11] desde ('localhost', 9266)
(debug) Recibido [DATA seq=12 data=b'304'] sin ACK previo, se asume el ACK perdido
(debug) Recibido [DATA seq=12 data=b'304'], enviando ACK seq=15
\end{lstlisting}

En otra ejecución, el servidor alcanzó a retransmitir el \texttt{SYN+ACK} antes de que
llegaran datos. El cliente, que ya estaba en \texttt{send}, lo reconoció como repetido y
reenvió el \texttt{ACK} del handshake, con lo que el servidor terminó \texttt{accept}:

\begin{lstlisting}
# cliente
(debug) Recibido [SYN+ACK seq=51] desde ('127.0.0.1', 9356), enviando ACK seq=52
(debug) Enviando [DATA seq=52 data=b'16']
(debug) Recibido [SYN+ACK seq=51] repetido, reenviando ACK del handshake
# servidor
(debug) Enviando [SYN+ACK seq=51] desde ('localhost', 9356)
(debug) Timeout esperando ACK del handshake, retransmitiendo SYN+ACK
(debug) Enviando [SYN+ACK seq=51] desde ('localhost', 9356)
(debug) Recibido [ACK seq=52], handshake completo
\end{lstlisting}

\subsection{Paso 8: cierre con pérdidas}

\begin{itemize}[nosep]
  \item \texttt{close()}: si no llega el \texttt{FIN+ACK} antes del timeout, reenvía el
        \texttt{FIN} y renueva el timeout, hasta 3 timeouts. Si se cumple el tercero,
        asume que la contraparte cerró y cierra el socket. Si recibe el \texttt{FIN+ACK},
        envía el último \texttt{ACK} 3 veces, esperando un timeout entre envíos.
  \item \texttt{recv\_close()}: después de enviar el \texttt{FIN+ACK} espera el último
        \texttt{ACK}; si no llega, reenvía el \texttt{FIN+ACK}, hasta 3 timeouts. Después
        del tercero asume que la contraparte cerró. Mientras espera el \texttt{FIN},
        también reenvía el \texttt{ACK} de cualquier segmento de datos duplicado, porque
        el emisor puede seguir retransmitiendo el último segmento si se perdió su
        \texttt{ACK}.
\end{itemize}

\textbf{Test.} En una ejecución con \texttt{netem} ninguno de los tres \texttt{FIN+ACK}
llegó al cliente; ambos lados terminaron igual de cerrar la conexión:

\begin{lstlisting}
# cliente (close)                            # servidor (recv_close)
(debug) Enviando [FIN seq=363]                (debug) Recibido [FIN seq=363]
(debug) Timeout 1/3 esperando FIN+ACK         (debug) Enviando [FIN+ACK seq=364]
(debug) Enviando [FIN seq=363]                (debug) Timeout 1/3 esperando el último ACK
(debug) Timeout 2/3 esperando FIN+ACK         (debug) Enviando [FIN+ACK seq=364]
(debug) Enviando [FIN seq=363]                (debug) Timeout 2/3 esperando el último ACK
(debug) Timeout 3/3 esperando FIN+ACK         (debug) Enviando [FIN+ACK seq=364]
(debug) No llegó FIN+ACK, se asume            (debug) Timeout 3/3 esperando el último ACK
        que la contraparte cerró              (debug) No llegó el último ACK, se asume
                                                      que la contraparte cerró
\end{lstlisting}

\section{Pruebas}

\subsection{Pérdidas con \texttt{netem}}
\label{sec:netem}

Las pruebas con pérdida se hicieron con \texttt{netem}, con la configuración del
enunciado:

\begin{lstlisting}
% sudo tc qdisc add dev lo root netem loss 20.0% delay 0.5s
% python3 servidor.py localhost 8000 > recibido.txt
% python3 cliente.py localhost 8000 < archivo.txt
% cmp archivo.txt recibido.txt
% sudo tc qdisc del dev lo root netem
\end{lstlisting}

Con el delay de 0,5 s en cada sentido, el RTT en \texttt{localhost} es de 1 s (medido con
\texttt{ping}), menor que el timeout de 1,5 s. Cada prueba se repitió 15 veces:

\begin{center}
\begin{tabular}{@{}lcc@{}}
  \toprule
  Prueba (\texttt{netem}: 20\% de pérdida, delay 0,5 s) & Ejecuciones exitosas & Criterio \\
  \midrule
  \texttt{cliente.py} / \texttt{servidor.py} con \texttt{archivo.txt} & 15/15 & \texttt{cmp} idéntico y ambos terminan \\
  \texttt{test\_cliente.py} / \texttt{test\_servidor.py} & 15/15 & 4 de 4 tests \texttt{Passed} \\
  \bottomrule
\end{tabular}
\end{center}

Cada envío del archivo tomó entre 20 y 43 segundos. En las 30 ejecuciones hubo 257
timeouts esperando ACKs de datos y el receptor recibió 120 segmentos duplicados. En el
handshake, el cliente retransmitió el \texttt{SYN} 15 veces y el servidor el
\texttt{SYN+ACK} 13 veces. El ACK implícito se usó 4 veces, el reenvío del \texttt{ACK}
ante un \texttt{SYN+ACK} repetido 3 veces, y los cierres por 3 timeouts ocurrieron 2
veces en cada lado.

\subsection{Envío de un archivo con pérdidas}

Este extracto del debug de una de las ejecuciones con \texttt{netem} muestra los dos
tipos de pérdida de Stop \& Wait. \texttt{netem} no informa qué segmento se perdió, pero
se deduce del lado del servidor. El segmento \texttt{seq=16} llegó una sola vez, así que
se perdió su primera transmisión y el emisor lo retransmitió al cumplirse el timeout. El
segmento \texttt{seq=64} llegó dos veces, así que lo que se perdió fue su
\texttt{ACK}: el receptor descartó el duplicado y reenvió el \texttt{ACK}.

\begin{lstlisting}
# cliente
(debug) Enviando [DATA seq=16 data=b'Redes CC4303 - A']
(debug) Timeout esperando ACK seq=32, retransmitiendo
(debug) Enviando [DATA seq=16 data=b'Redes CC4303 - A']
...
(debug) Enviando [DATA seq=64 data=b'vo tiene bastant']
(debug) Timeout esperando ACK seq=80, retransmitiendo
(debug) Enviando [DATA seq=64 data=b'vo tiene bastant']
# servidor
(debug) Recibido [DATA seq=16 data=b'Redes CC4303 - A'], enviando ACK seq=32
...
(debug) Recibido [DATA seq=64 data=b'vo tiene bastant'], enviando ACK seq=80
(debug) Recibido [DATA seq=64 data=b'vo tiene bastant'] duplicado, reenviando ACK seq=80
\end{lstlisting}

\subsection{\texttt{recv} con \texttt{buff\_size < message\_length}}

A los tests del enunciado se agregó el test 4: se envía un mensaje de $2n = 42$ bytes y
se llama dos veces a \texttt{recv(n)} con $n = 21$, que no es múltiplo de 16. La primera
llamada recibe dos segmentos (32 bytes), retorna 21 y deja 11 en \texttt{recv\_buffer}.
La segunda usa esos 11 bytes y recibe el último segmento de 10. El test 3 cubre el mismo
caso con $n = 14$.

\begin{lstlisting}
Test 1 received: b'Mensje de len=16'
Test 1: Passed
Test 2 received: b'Mensaje de largo 19'
Test 2: Passed
Test 3 received: b'Mensaje de largo 19'
Test 3: Passed
Test 4 received: b'Primera mitad de 21 b' b'Segunda mitad de 21 b'
Test 4: Passed
Conexión cerrada
\end{lstlisting}

\section{Casos borde y limitaciones}

\begin{itemize}
  \item \textbf{Se pierden los 3 \texttt{FIN}.} \texttt{close()} asume que la contraparte
        cerró, pero \texttt{recv\_close()} espera el \texttt{FIN} sin timeout y queda
        bloqueada. El enunciado solo exige que el cierre funcione siempre desde el lado
        que llama a \texttt{close()} y, en la contraparte, una vez que recibe un
        \texttt{FIN}. Con el 20\% de pérdida de \texttt{netem}, la probabilidad de este caso
        es $0{,}2^3 = 0{,}8\%$ por cierre. No ocurrió en las pruebas con \texttt{netem}.
  \item \textbf{Se pierden los 3 últimos \texttt{ACK} del cierre.} \texttt{recv\_close()}
        termina igual después de 3 timeouts.
  \item \textbf{Se pierde el último \texttt{ACK} de datos antes del cierre.} El emisor
        retransmite el segmento mientras el receptor ya está en \texttt{recv\_close()};
        este reenvía el \texttt{ACK} y el emisor sigue normalmente con \texttt{close()}.
  \item \textbf{\texttt{SYN} retransmitidos.} Quedan en el socket que escucha y se
        ignoran porque la conexión ya fue aceptada.
  \item \textbf{Handshake y envío de datos sin límite de reintentos.} \texttt{connect},
        \texttt{accept} y \texttt{send} retransmiten indefinidamente. Si la contraparte
        deja de existir, quedan esperando para siempre.
  \item \textbf{El timeout se reinicia con cada segmento recibido.} Con
        \texttt{settimeout}, el plazo vuelve a empezar cada vez que \texttt{recvfrom}
        recibe algo, aunque no sea el \texttt{ACK} esperado. Si llegaran segmentos
        inesperados de forma continua, la retransmisión se atrasaría. En la práctica solo
        llegan algunos duplicados, por lo que el efecto es despreciable.
  \item \textbf{Timeout fijo.} El timeout no se adapta al RTT. Si el RTT supera 1,5
        segundos, habría retransmisiones innecesarias. Aun así la transferencia sería
        correcta, porque los duplicados se descartan por número de secuencia.
  \item \textbf{\texttt{buff\_size} del servidor.} \texttt{servidor.py} llama una vez a
        \texttt{recv} con un \texttt{buff\_size} de 1~MiB, por lo que recibe el archivo
        completo si mide menos que eso.
  \item \textbf{Mensaje de largo 0.} \texttt{send(b"")} envía solo el segmento con el
        largo \texttt{b'0'} y \texttt{recv} retorna \texttt{b""}.
\end{itemize}

\section*{Declaración de uso de IA}

En este informe se utilizó el modelo de IA Claude exclusivamente como apoyo en la
revisión del código, recomendaciones de diseño, la redacción del informe, la creación de
las tablas y la elaboración de los gráficos y diagramas en \LaTeX{}. Las decisiones de
diseño, la implementación del código, la ejecución de los experimentos y la obtención de
los datos fueron realizadas por el autor. Todo el contenido generado con apoyo de la
herramienta fue revisado y validado, y su responsabilidad es del autor.

\end{document}
